\documentclass[11pt,a4paper]{article}

\usepackage[T1]{fontenc}
\usepackage[utf8]{inputenc}
\usepackage[french]{babel}
\usepackage[a4paper,margin=2.2cm]{geometry}
\usepackage{amsmath,amssymb,amsfonts,amsthm}
\usepackage{mathtools}
\usepackage{bm}
\usepackage{graphicx}
\usepackage{xcolor}
\usepackage{hyperref}
\usepackage{enumitem}
\usepackage{booktabs}
\usepackage{array}
\usepackage{tikz}
\usepackage{fancyhdr}
\usepackage{caption}
\usepackage{subcaption}

\hypersetup{
    colorlinks=true,
    linkcolor=blue!60!black,
    urlcolor=blue!60!black,
    citecolor=blue!60!black,
    pdftitle={Grad-CAM -- Fondements mathematiques et interpretation},
    pdfauthor={Yann Sofian Smatti}
}

\setlength{\parindent}{0pt}
\setlength{\parskip}{0.6em}

\pagestyle{fancy}
\fancyhf{}
\lhead{Grad-CAM -- document pedagogique}
\rhead{\thepage}

\title{\textbf{Grad-CAM}\\
Fondements mathématiques, intuition visuelle et interprétation\\
\large Document détaillé pour comprendre l'explicabilité locale en vision par ordinateur}
\author{Yann Sofian Smatti}
\date{Mars 2026}

\begin{document}

\maketitle

\begin{center}
\textbf{Explainable AI · CNN · Imagerie médicale · Vision par ordinateur}
\end{center}

\tableofcontents
\newpage

\section{Introduction}

Les réseaux de neurones convolutifs obtiennent d'excellentes performances en classification d'images, mais ils sont souvent perçus comme des \emph{boîtes noires}. Un modèle peut prédire qu'une image contient une pathologie avec une forte confiance sans indiquer clairement \emph{pourquoi} il prend cette décision.

Cette absence d'interprétabilité est particulièrement problématique en imagerie médicale. Un clinicien ne peut pas se contenter d'une probabilité ; il veut savoir quelles régions de l'image ont influencé le modèle.

Grad-CAM, pour \emph{Gradient-weighted Class Activation Mapping}, est une méthode d'explicabilité locale qui répond à la question suivante :

\begin{center}
\emph{Quelles zones de l'image ont le plus contribué à la prédiction d'une classe donnée ?}
\end{center}

Le but de ce document est d'expliquer Grad-CAM :
\begin{itemize}
    \item intuitivement ;
    \item mathématiquement ;
    \item algorithmiquement ;
    \item avec un angle spécifique à l'imagerie médicale.
\end{itemize}

\section{Pourquoi a-t-on besoin d'explicabilité ?}

\subsection{Le problème de la boîte noire}

Considérons un réseau de neurones convolutif réalisant une classification binaire. À partir d'une image $x$, il produit

\[
P(\text{maladie} \mid x)=0.93.
\]

Cette sortie est utile, mais elle ne nous dit pas si le modèle regarde :
\begin{itemize}
    \item la véritable lésion ;
    \item un artefact d'acquisition ;
    \item une annotation textuelle ;
    \item un bord noir de l'image.
\end{itemize}

Deux modèles peuvent produire la même sortie numérique, mais pour des raisons très différentes.

\subsection{Pourquoi c'est critique en médecine}

En imagerie médicale, une bonne explication visuelle permet de :
\begin{itemize}
    \item vérifier que le réseau regarde la bonne région anatomique ;
    \item détecter des biais du dataset ;
    \item repérer des comportements aberrants ;
    \item renforcer ou diminuer la confiance dans le modèle.
\end{itemize}

Grad-CAM ne remplace pas une validation clinique, mais fournit une \textbf{explication locale} de la décision.

\section{Ce qu'un CNN voit réellement}

\subsection{L'image comme matrice}

Une image en niveaux de gris est une matrice

\[
I \in \mathbb{R}^{H \times W},
\]

et une image couleur un tenseur

\[
I \in \mathbb{R}^{H \times W \times 3}.
\]

Le réseau ne voit pas une \og lésion \fg{} ou un \og poumon \fg{}. Il voit seulement des nombres.

\subsection{Convolutions et cartes de caractéristiques}

Dans un CNN, les premières couches appliquent des filtres convolutifs qui extraient des motifs locaux :
\begin{itemize}
    \item bords ;
    \item textures ;
    \item petits contrastes ;
    \item structures plus complexes à mesure que l'on monte dans le réseau.
\end{itemize}

Chaque filtre produit une \textbf{carte de caractéristiques}, aussi appelée \emph{feature map}. On les note généralement

\[
A^1, A^2, \dots, A^K.
\]

Chaque carte $A^k$ est une image intermédiaire, souvent de plus petite taille spatiale que l'image d'origine.

\subsection{Interprétation intuitive}

On peut voir une feature map comme un détecteur spatial de motifs. Si une région donnée de l'image active fortement un certain filtre, alors la valeur correspondante dans la feature map sera grande.

Dans les couches profondes, ces cartes ne détectent plus seulement des bords simples mais des motifs plus abstraits et plus utiles pour la classe finale.

\section{Objectif de Grad-CAM}

Grad-CAM cherche à répondre à cette question :

\begin{center}
\emph{Parmi les cartes de caractéristiques finales, lesquelles sont les plus importantes pour la classe cible, et où sont-elles actives spatialement ?}
\end{center}

L'idée générale est de :
\begin{enumerate}[label=\arabic*)]
    \item choisir une classe cible $c$ ;
    \item mesurer l'influence de chaque feature map sur le score de cette classe ;
    \item combiner les feature maps avec ces poids d'influence ;
    \item obtenir une heatmap superposable à l'image originale.
\end{enumerate}

\section{De la prédiction à la dérivée}

\subsection{Score de classe}

Soit $y^c$ le score associé à la classe $c$. Selon l'implémentation, $y^c$ peut désigner :
\begin{itemize}
    \item le logit avant softmax ;
    \item ou plus rarement la probabilité après softmax.
\end{itemize}

En pratique, on préfère souvent utiliser le \textbf{logit} pour éviter certains effets de saturation.

\subsection{Pourquoi regarder le gradient ?}

Le gradient permet de mesurer la sensibilité du score $y^c$ aux activations internes du réseau. Si l'on note par $A_{ij}^k$ la valeur de la feature map $k$ à la position spatiale $(i,j)$, alors

\[
\frac{\partial y^c}{\partial A_{ij}^k}
\]

mesure l'effet d'une petite variation de $A_{ij}^k$ sur le score de la classe $c$.

\textbf{Interprétation :}
\begin{itemize}
    \item si cette dérivée est grande et positive, augmenter cette activation augmente le score de la classe ;
    \item si elle est négative, cette activation pousse plutôt contre la classe ;
    \item si elle est proche de zéro, cette activation a peu d'influence.
\end{itemize}

\section{Formulation mathématique de Grad-CAM}

\subsection{Poids d'importance par canal}

Pour chaque canal $k$ de la dernière couche convolutive, Grad-CAM définit un poids

\[
\alpha_k^c = \frac{1}{Z}\sum_i\sum_j \frac{\partial y^c}{\partial A_{ij}^k},
\]

où $Z$ est le nombre total de positions spatiales :

\[
Z = H'W'.
\]

Ici, $H'$ et $W'$ sont la hauteur et la largeur des feature maps de la dernière couche convolutionnelle.

\subsection{Sens de la moyenne spatiale}

Le terme

\[
\frac{1}{Z}\sum_i\sum_j \frac{\partial y^c}{\partial A_{ij}^k}
\]

est une moyenne globale des gradients du canal $k$.

\textbf{Idée :}
au lieu d'examiner séparément chaque position, on résume l'influence globale du canal $k$ sur la classe $c$ par un seul scalaire $\alpha_k^c$.

Ainsi :
\begin{itemize}
    \item un canal très utile pour la classe obtient un poids élevé ;
    \item un canal peu utile obtient un poids faible ;
    \item un canal nuisible peut obtenir un poids négatif.
\end{itemize}

\subsection{Combinaison linéaire des cartes}

Une fois les poids calculés, la carte Grad-CAM est définie par

\[
L_{\mathrm{GradCAM}}^c = \mathrm{ReLU}\left(\sum_k \alpha_k^c A^k\right).
\]

Il s'agit d'une combinaison linéaire pondérée des feature maps.

\subsection{Pourquoi la fonction ReLU ?}

La fonction

\[
\mathrm{ReLU}(x)=\max(0,x)
\]

supprime les contributions négatives.

L'idée est la suivante : pour expliquer \emph{pourquoi} le modèle a choisi la classe $c$, on préfère mettre en évidence les régions qui \textbf{augmentent} le score de cette classe, plutôt que celles qui le diminuent.

Autrement dit, Grad-CAM cherche les \emph{preuves positives} en faveur de la classe cible.

\section{Exemple jouet détaillé}

Considérons trois feature maps finales :

\[
A^1=
\begin{pmatrix}
0 & 1 & 0 \\
1 & 3 & 1 \\
0 & 1 & 0
\end{pmatrix},
\qquad
A^2=
\begin{pmatrix}
0 & 0 & 1 \\
0 & 2 & 2 \\
1 & 2 & 1
\end{pmatrix},
\qquad
A^3=
\begin{pmatrix}
1 & 2 & 1 \\
2 & 8 & 2 \\
1 & 2 & 1
\end{pmatrix}.
\]

Supposons que pour la classe cible $c$, les poids obtenus soient :

\[
\alpha_1^c = 0.1,
\qquad
\alpha_2^c = 0.3,
\qquad
\alpha_3^c = 0.9.
\]

Alors,

\[
L^c = \mathrm{ReLU}(0.1A^1 + 0.3A^2 + 0.9A^3).
\]

Calculons :

\[
0.1A^1=
\begin{pmatrix}
0 & 0.1 & 0 \\
0.1 & 0.3 & 0.1 \\
0 & 0.1 & 0
\end{pmatrix},
\]

\[
0.3A^2=
\begin{pmatrix}
0 & 0 & 0.3 \\
0 & 0.6 & 0.6 \\
0.3 & 0.6 & 0.3
\end{pmatrix},
\]

\[
0.9A^3=
\begin{pmatrix}
0.9 & 1.8 & 0.9 \\
1.8 & 7.2 & 1.8 \\
0.9 & 1.8 & 0.9
\end{pmatrix}.
\]

La somme vaut

\[
L^c_{\text{avant ReLU}}=
\begin{pmatrix}
0.9 & 1.9 & 1.2 \\
1.9 & 8.1 & 2.5 \\
1.2 & 2.5 & 1.2
\end{pmatrix}.
\]

Toutes les valeurs étant positives, ReLU ne change rien ici.

\textbf{Interprétation :}
la plus forte valeur se situe au centre. La heatmap finale indiquera donc que le modèle regarde essentiellement la région centrale pour justifier sa prédiction.

\section{Pourquoi la dernière couche convolutive ?}

Grad-CAM s'applique généralement à la \textbf{dernière couche convolutionnelle}.

Pourquoi ? Parce qu'elle offre un compromis entre :
\begin{itemize}
    \item \textbf{richesse sémantique} : les couches profondes représentent des motifs complexes utiles à la classification ;
    \item \textbf{structure spatiale} : contrairement aux couches denses finales, les convolutions conservent une organisation dans l'espace.
\end{itemize}

Les couches très basses ont une grande résolution spatiale, mais leurs activations sont souvent trop primitives. Les couches totalement connectées contiennent une information très abstraite, mais perdent la localisation spatiale.

La dernière couche convolutionnelle est donc l'endroit naturel pour produire une carte explicative locale.

\section{Algorithme de Grad-CAM}

\subsection{Étapes conceptuelles}

Soit un réseau convolutif et une image d'entrée $x$.

\begin{enumerate}[label=\arabic*)]
    \item Faire une passe avant dans le réseau pour obtenir le score $y^c$ de la classe cible.
    \item Extraire les feature maps $A^k$ de la dernière couche convolutionnelle.
    \item Calculer les gradients de $y^c$ par rapport à chaque élément $A_{ij}^k$.
    \item Calculer les poids $\alpha_k^c$ par moyenne spatiale des gradients.
    \item Former la combinaison pondérée $\sum_k \alpha_k^c A^k$.
    \item Appliquer ReLU.
    \item Redimensionner la heatmap à la taille de l'image originale.
    \item Superposer la heatmap à l'image.
\end{enumerate}

\subsection{Pseudo-code}

\begin{center}
\begin{minipage}{0.92\linewidth}
\hrule
\vspace{0.4em}
\textbf{Algorithme Grad-CAM}
\vspace{0.4em}
\hrule
\vspace{0.6em}
Entrée : image $x$, modèle $f$, classe cible $c$ \\
Sortie : heatmap explicative $L^c$

\begin{enumerate}[label=\arabic*)]
    \item Calculer la passe avant et stocker les activations $A^k$.
    \item Calculer le score de classe $y^c$.
    \item Calculer les gradients $\partial y^c / \partial A_{ij}^k$.
    \item Pour chaque canal $k$, calculer
    \[
    \alpha_k^c = \frac{1}{Z}\sum_{i,j} \frac{\partial y^c}{\partial A_{ij}^k}.
    \]
    \item Combiner les cartes :
    \[
    L^c = \mathrm{ReLU}\left(\sum_k \alpha_k^c A^k\right).
    \]
    \item Normaliser et interpoler la carte à la taille de l'image d'origine.
\end{enumerate}
\vspace{0.6em}
\hrule
\end{minipage}
\end{center}

\section{Interprétation géométrique}

On peut voir les activations locales de la dernière couche convolutive comme un ensemble de directions ou de motifs dans un espace de représentation. Le score de classe $y^c$ dépend de ces activations.

Le gradient indique comment ce score varie localement. Les poids $\alpha_k^c$ résument l'importance de chaque canal dans cette variation.

Ainsi, Grad-CAM réalise deux opérations :
\begin{itemize}
    \item une \textbf{pondération sémantique} des canaux par les gradients ;
    \item une \textbf{projection spatiale} des activations utiles à la classe.
\end{itemize}

On peut résumer cela par :

\begin{center}
\emph{Heatmap = où le réseau s'active $\times$ à quel point cette activation aide la classe cible.}
\end{center}

\section{Superposition à l'image originale}

La carte $L_{\mathrm{GradCAM}}^c$ est de petite taille, par exemple $7 \times 7$ ou $14 \times 14$. Pour être lisible, on :
\begin{itemize}
    \item la normalise dans $[0,1]$ ;
    \item l'interpole à la taille de l'image d'entrée ;
    \item la représente sous forme de heatmap ;
    \item la superpose à l'image originale.
\end{itemize}

Mathématiquement, si $\widetilde{L}^c$ est la version interpolée, on peut construire une image explicative de type

\[
I_{\mathrm{exp}} = (1-\lambda)I + \lambda H(\widetilde{L}^c),
\]

où $H$ désigne une colorisation thermique et $\lambda \in [0,1]$ un coefficient de mélange.

\section{Exemple de lecture d'une heatmap}

Supposons qu'un modèle de radiographie thoracique prédise :

\[
P(\text{pneumonie})=0.91.
\]

Trois situations sont possibles :

\begin{enumerate}[label=\arabic*)]
    \item La heatmap se concentre sur une opacité pulmonaire cohérente avec la pathologie : c'est rassurant.
    \item La heatmap se concentre sur le bord de l'image ou sur une annotation : cela révèle un biais.
    \item La heatmap est diffuse, peu localisée, ou instable : cela suggère une explication peu robuste.
\end{enumerate}

Grad-CAM n'est donc pas seulement un outil de visualisation esthétique ; c'est aussi un outil de \textbf{diagnostic du modèle lui-même}.

\section{Spécificités en imagerie médicale}

\subsection{Points forts}

En imagerie médicale, Grad-CAM est utile pour :
\begin{itemize}
    \item vérifier la cohérence anatomique ;
    \item produire des visualisations compréhensibles par les cliniciens ;
    \item détecter des corrélations parasites ;
    \item documenter qualitativement les sorties du modèle.
\end{itemize}

\subsection{Limites importantes}

Il faut cependant rester prudent.

Grad-CAM :
\begin{itemize}
    \item n'est pas une preuve causale ;
    \item dépend du choix de la couche convolutionnelle ;
    \item peut produire des heatmaps visuellement séduisantes mais peu informatives ;
    \item donne une explication \textbf{locale} pour une prédiction donnée, pas une compréhension globale du modèle.
\end{itemize}

\section{Grad-CAM, CAM et variantes}

\subsection{CAM classique}

La méthode CAM originale suppose une architecture particulière avec :
\begin{itemize}
    \item une dernière couche convolutionnelle ;
    \item un global average pooling ;
    \item une couche linéaire finale.
\end{itemize}

Grad-CAM généralise cette idée en utilisant les gradients, ce qui lui permet de fonctionner avec une plus grande variété d'architectures.

\subsection{Grad-CAM++}

Grad-CAM++ modifie la pondération des activations pour mieux gérer les cas où plusieurs occurrences d'un même objet apparaissent dans l'image.

\subsection{Guided Grad-CAM}

Guided Grad-CAM combine Grad-CAM avec des méthodes de backpropagation guidée afin d'obtenir des cartes plus fines, mais souvent moins directement interprétables cliniquement.

\section{Pourquoi Grad-CAM marche-t-il souvent bien ?}

Deux raisons principales expliquent son succès :
\begin{enumerate}[label=\arabic*)]
    \item Les dernières convolutions conservent une structure spatiale.
    \item Les gradients donnent une mesure locale de l'influence sur la classe.
\end{enumerate}

Le mélange des deux fournit une explication à la fois :
\begin{itemize}
    \item sémantique ;
    \item localisée ;
    \item relativement simple à calculer.
\end{itemize}

\section{Exemple conceptuel de pipeline}

\begin{center}
\begin{tikzpicture}[scale=0.95, every node/.style={font=\small}]
    \draw[rounded corners, thick] (0,0) rectangle (2.5,1.2);
    \node at (1.25,0.6) {Image d'entrée};

    \draw[->, thick] (2.5,0.6) -- (4,0.6);

    \draw[rounded corners, thick] (4,0) rectangle (6.8,1.2);
    \node at (5.4,0.6) {CNN};

    \draw[->, thick] (6.8,0.6) -- (8.5,0.6);

    \draw[rounded corners, thick] (8.5,0) rectangle (11.8,1.2);
    \node at (10.15,0.6) {Feature maps $A^k$};

    \draw[->, thick] (10.15,0) -- (10.15,-1.4);
    \node at (11.9,-0.7) {Gradients $\partial y^c / \partial A^k$};

    \draw[->, thick] (11.8,0.6) -- (13.4,0.6);

    \draw[rounded corners, thick] (13.4,0) rectangle (15.9,1.2);
    \node at (14.65,0.6) {Grad-CAM};

    \draw[->, thick] (15.9,0.6) -- (17.5,0.6);

    \draw[rounded corners, thick] (17.5,0) rectangle (20,1.2);
    \node at (18.75,0.6) {Heatmap};
\end{tikzpicture}
\end{center}

\section{Implémentation pratique}

En pratique, avec TensorFlow ou PyTorch, on procède ainsi :
\begin{itemize}
    \item on construit un modèle auxiliaire qui renvoie à la fois les activations de la dernière couche convolutionnelle et la sortie du réseau ;
    \item on utilise l'autodifférentiation pour calculer les gradients ;
    \item on calcule les poids moyens par canal ;
    \item on forme la heatmap ;
    \item on l'affiche après interpolation.
\end{itemize}

Cette méthode est rapide, ne nécessite pas de réentraîner le modèle et s'applique image par image.

\section{Points de vigilance méthodologiques}

Quand on utilise Grad-CAM dans un projet sérieux, il faut :
\begin{itemize}
    \item vérifier plusieurs images, pas seulement un exemple réussi ;
    \item comparer des cas vrais positifs, faux positifs et faux négatifs ;
    \item inspecter visuellement si le modèle regarde la région pertinente ;
    \item conserver un esprit critique sur l'interprétation.
\end{itemize}

Une heatmap convaincante sur trois images n'implique pas que le modèle est fiable globalement.

\section{Résumé intuitif}

On peut résumer Grad-CAM de la manière suivante :

\begin{enumerate}[label=\arabic*)]
    \item Le CNN transforme l'image en cartes internes de motifs.
    \item Certaines de ces cartes aident davantage la classe cible que d'autres.
    \item Les gradients permettent de mesurer cette aide.
    \item On recombine ces cartes avec leurs poids d'importance.
    \item On obtient une carte de chaleur qui indique où le modèle regarde.
\end{enumerate}

\begin{center}
\fbox{\parbox{0.9\linewidth}{
\centering
\textbf{Grad-CAM montre quelles régions de l'image contribuent positivement à la prédiction d'une classe donnée.}
}}
\end{center}

\section{Comparaison avec d'autres méthodes d'explicabilité}

\subsection{Class Activation Mapping (CAM)}

La méthode CAM (Class Activation Mapping) est l'ancêtre direct de Grad-CAM.

Elle suppose une architecture spécifique :

\begin{itemize}
\item dernière couche convolutionnelle
\item global average pooling
\item couche linéaire finale
\end{itemize}

Dans ce cas, la carte d'activation de classe peut être calculée directement :

\[
L^c = \sum_k w_k^c A^k
\]

où

\begin{itemize}
\item $A^k$ est la feature map $k$
\item $w_k^c$ est le poids de la classe $c$ dans la couche finale
\end{itemize}

Cette approche est simple mais impose une architecture particulière.

Grad-CAM généralise cette idée en utilisant les gradients.

\subsection{Grad-CAM}

Grad-CAM fonctionne avec n'importe quel CNN.

La pondération est obtenue via les gradients :

\[
\alpha_k^c =
\frac{1}{Z}
\sum_{i,j}
\frac{\partial y^c}{\partial A_{ij}^k}
\]

et la carte finale :

\[
L^c =
ReLU
\left(
\sum_k
\alpha_k^c A^k
\right)
\]

Cette méthode permet donc d'expliquer les prédictions sans modifier l'architecture du réseau.

\subsection{Grad-CAM++}

Grad-CAM++ améliore la localisation lorsque plusieurs objets sont présents.

Les poids deviennent :

\[
\alpha_{ij}^{kc}
\]

et sont calculés avec des dérivées d'ordre supérieur.

Cela permet :

\begin{itemize}
\item meilleure localisation
\item détection de multiples objets
\end{itemize}

mais augmente le coût de calcul.

\section{Implémentation TensorFlow}

Une implémentation simplifiée est la suivante.

\begin{verbatim}
with tf.GradientTape() as tape:

    conv_outputs, predictions = grad_model(image)

    loss = predictions[:, class_index]

grads = tape.gradient(loss, conv_outputs)

pooled_grads = tf.reduce_mean(grads, axis=(0,1,2))

conv_outputs = conv_outputs[0]

heatmap = tf.reduce_sum(pooled_grads * conv_outputs, axis=-1)

heatmap = tf.maximum(heatmap,0)

heatmap = heatmap / tf.reduce_max(heatmap)
\end{verbatim}

Les étapes correspondent exactement à l'algorithme mathématique.

\section{Application dans le projet MedVision}

Dans le projet MedVision, Grad-CAM est utilisé pour interpréter les prédictions du modèle CNN entraîné sur des radiographies.

Le pipeline complet est :

\begin{enumerate}

\item Chargement de l'image médicale

\item Prétraitement et normalisation

\item Prédiction par le CNN

\item Extraction de la dernière couche convolutionnelle

\item Calcul de la heatmap Grad-CAM

\item Redimensionnement

\item Superposition avec l'image originale

\end{enumerate}

Cette heatmap permet de vérifier si le modèle regarde réellement une région pathologique.

\section{Analyse d'une heatmap médicale}

Trois cas sont possibles.

\textbf{Cas 1 : heatmap cohérente}

La heatmap est localisée sur la lésion.

Cela suggère que le modèle a appris un signal clinique pertinent.

\textbf{Cas 2 : heatmap parasite}

La heatmap est localisée sur :

\begin{itemize}

\item un bord d'image
\item une annotation
\item un artefact du scanner

\end{itemize}

Cela révèle un biais du dataset.

\textbf{Cas 3 : heatmap diffuse}

La carte est dispersée et peu informative.

Cela peut indiquer :

\begin{itemize}

\item un modèle instable
\item un manque de données
\item une mauvaise architecture

\end{itemize}

\section{Bonnes pratiques}

Pour analyser un modèle avec Grad-CAM :

\begin{itemize}

\item inspecter plusieurs images

\item analyser vrais positifs et faux positifs

\item vérifier la cohérence anatomique

\item comparer différentes couches convolutionnelles

\end{itemize}

Grad-CAM est un outil d'analyse du modèle, pas une preuve absolue.


\section{Conclusion}

Grad-CAM est une méthode d'explicabilité locale élégante, simple et très utilisée en vision par ordinateur. Elle combine deux idées fondamentales :
\begin{itemize}
    \item la structure spatiale conservée par les couches convolutionnelles ;
    \item l'information de sensibilité fournie par les gradients.
\end{itemize}

Son succès vient du fait qu'elle offre une explication visuelle directe, souvent suffisamment parlante pour être utilisée dans des contextes appliqués, notamment en imagerie médicale.

Il faut toutefois garder en tête qu'il s'agit d'un outil d'interprétation \emph{partiel}. Grad-CAM aide à comprendre une décision, mais ne résout pas à lui seul toutes les questions de robustesse, de causalité ou de confiance.

\end{document}
```
